% Part: propositional-logic
% Chapter: syntax-and-semantics
% Section: formulas

\documentclass[../../../include/open-logic-section]{subfiles}

\begin{document}

\olfileid{pl}{syn}{fml}

\olsection{\usetoken{P}{formula}ی گزاره‌ای}

!!^{formula}sی منطق گزاره‌ای از
\emph{!!{propositional
    variable}s}\iftag{prvFalse}{\iftag{prvTrue}{،}{ و} ثابت
  گزاره‌ایِ~$\lfalse$}{}\iftag{prvTrue}{\iftag{prvFalse}{ و}{ و} ثابت
  گزاره‌ایِ~$\ltrue$}{} با استفاده از \emph{رابط‌های منطقی} ساخته می‌شوند.

\begin{enumerate}
\item یک مجموعهٔ !!^a{denumerable} به نام~$\PVar$ از
  !!{propositional variable}s، یعنی $\Obj p_0$،
  $\Obj p_1$، \dots
\tagitem{prvFalse}{ثابت گزاره‌ایِ !!{falsity}~$\lfalse$.}{}
\tagitem{prvTrue}{ثابت گزاره‌ایِ !!{truth}~$\ltrue$.}{}
\item رابط‌های منطقی:
  \startycommalist
  \iftag{prvNot}{\ycomma $\lnot$ (نفی)}{}%
  \iftag{prvAnd}{\ycomma $\land$ (عطف)}{}%
  \iftag{prvOr}{\ycomma $\lor$ (فصل)}{}%
  \iftag{prvIf}{\ycomma $\lif$ (!!{conditional})}{}%
  \iftag{prvIff}{\ycomma $\liff$ (!!{biconditional})}{}%
\item نشانه‌های سجاوندی: ( و )، و نیز ویرگول.
\end{enumerate}

این زبان منطق گزاره‌ای را با $\Lang L_0$ نشان می‌دهیم.

\iftag{defNot,defOr,defAnd,defIf,defIff,defTrue,defFalse,defEx,defAll}{%
افزون بر رابط‌های اولیه‌ای که در بالا معرفی شدند، از نمادهای
\emph{تعریف‌شدهٔ} زیر نیز استفاده می‌کنیم:
\startycommalist
  \iftag{defNot}{\ycomma $\lnot$ (نفی)}{}%
  \iftag{defAnd}{\ycomma $\land$ (عطف)}{}%
  \iftag{defOr}{\ycomma $\lor$ (فصل)}{}%
  \iftag{defIf}{\ycomma $\lif$ (!!{conditional})}{}%
  \iftag{defIff}{\ycomma $\liff$ (!!{biconditional})}{}%
  \iftag{defFalse}{\ycomma $\lfalse$ (!!{falsity})}{}%
  \iftag{defTrue}{\ycomma $\ltrue$ (!!{truth})}}{}.

\begin{tagblock}{defNot,defOr,defAnd,defIf,defIff,defTrue,defFalse,defEx,defAll}
\begin{explain}
نماد تعریف‌شده رسماً جزئی از زبان نیست، بلکه به‌عنوان اختصاری غیرصوری
معرفی می‌شود: چنین نمادی به ما امکان می‌دهد فرمول‌هایی را کوتاه کنیم که
اگر فقط از نمادهای اولیه استفاده می‌کردیم، بسیار طولانی می‌شدند. این مزیت
آشکار است؛ اما مزیت بزرگ‌تر آن است که برهان‌ها کوتاه‌تر می‌شوند. اگر نمادی
اولیه باشد، باید در برهان‌ها جداگانه بررسی شود. بنابراین هرچه شمار نمادهای
اولیه بیشتر باشد، برهان‌های ما طولانی‌تر خواهند بود.
\end{explain}
\end{tagblock}

% Alternate symbols

\begin{intro}
ممکن است با اصطلاح‌ها و نمادهایی متفاوت از آنچه در بالا به کار بردیم آشنا
باشید. متون منطق (و مدرسان) معمولاً یکی از نمادهای $\sim$، $\neg$ و~!\ را
برای «نفی»، و یکی از نمادهای $\wedge$، $\cdot$ و $\&$ را برای «عطف» به
کار می‌برند. نمادهای رایج برای «شرطی» یا «استلزام» عبارت‌اند از
$\rightarrow$، $\Rightarrow$ و $\supset$.
\iftag{prvIff,defIff}{نمادهای «دوشرطی»، «استلزام دوسویه» یا «هم‌ارزی
  (مادی)» عبارت‌اند از $\leftrightarrow$،
  $\Leftrightarrow$ و $\equiv$.}{}
\iftag{prvFalse,defFalse}{نماد $\lfalse$ را به نام‌های «کذب»، «فالسوم»،
  «محال» یا «پایین» نیز می‌خوانند.}{}
\iftag{prvTrue,defTrue}{نماد $\ltrue$ را به نام‌های «صدق»، «وروم» یا
  «بالا» نیز می‌خوانند.}{}
\end{intro}

\begin{defn}[فرمول]
\ollabel{defn:formulas}
مجموعهٔ~$\Frm[L_0]$ از \emph{!!{formula}sی} منطق گزاره‌ای به‌صورت استقرایی
چنین تعریف می‌شود:
\begin{enumerate}
\tagitem{prvFalse}{$\lfalse$ یک !!{formula} اتمی است.}{}

\tagitem{prvTrue}{$\ltrue$ یک !!{formula} اتمی است.}{}

\item هر !!{propositional variable}~$\Obj p_i$ یک !!{formula} اتمی است.

\tagitem{prvNot}{اگر $!A$ یک !!a{formula} باشد، آنگاه $\lnot !A$ یک
  !!a{formula} است.}{}

\tagitem{prvAnd}{اگر $!A$ و $!B$ از !!{formula}s باشند، آنگاه $(!A \land
  !B)$ یک !!a{formula} است.}{}

\tagitem{prvOr}{اگر $!A$ و $!B$ از !!{formula}s باشند، آنگاه $(!A \lor !B)$
  یک !!a{formula} است.}{}

\tagitem{prvIf}{اگر $!A$ و $!B$ از !!{formula}s باشند، آنگاه $(!A \lif !B)$
  یک !!a{formula} است.}{}

\tagitem{prvIff}{اگر $!A$ و $!B$ از !!{formula}s باشند، آنگاه $(!A \liff !B)$
  یک !!a{formula} است.}{}

\tagitem{limitClause}{هیچ چیز دیگری !!a{formula} نیست.}{}
\end{enumerate}
\end{defn}

\begin{explain}
تعریف !!{formula}s یک \emph{تعریف استقرایی} است. در اصل، مجموعهٔ
!!{formula}s را در بی‌نهایت مرحله می‌سازیم. در مرحلهٔ آغازین، همهٔ
فرمول‌های اتمی را فرمول اعلام می‌کنیم؛ این مرحله با چند حالت نخست تعریف،
یعنی حالت‌های
\iftag{prvTrue}{$\ltrue$، }{}%
\iftag{prvFalse}{$\lfalse$، }{}%
$\Obj p_i$ متناظر است. پس «!!{formula} اتمی» یعنی هر !!{formula}ی که یکی از
این صورت‌ها را داشته باشد.

حالت‌های دیگر تعریف، قواعدی برای ساختن !!{formula}sی تازه از
!!{formula}sی پیش‌ترساخته به دست می‌دهند. در مرحلهٔ دوم می‌توانیم
با آن قواعد !!{formula}s را از !!{formula}sی اتمی بسازیم. در
مرحلهٔ سوم، فرمول‌های تازه را از فرمول‌های اتمی و فرمول‌های به‌دست‌آمده در
مرحلهٔ دوم می‌سازیم و همین‌طور ادامه می‌دهیم. اصطلاح «!!{formula}» بر هر
چیزی اطلاق می‌شود که سرانجام در یکی از این مراحل ساخته شود، و نه هیچ چیز
دیگر.
\end{explain}

هنگام نوشتن فرمول $(!B \ast !C)$ که از $!B$ و $!C$ با استفاده از یک رابط
دوجایگاهی~$\ast$ ساخته شده است، اغلب بیرونی‌ترین جفت پرانتز را حذف می‌کنیم
و فقط~$!B \ast !C$ می‌نویسیم.

\begin{tagblock}{defNot,defOr,defAnd,defIf,defIff,defTrue,defFalse,defEx,defAll}
\begin{defn}
فرمول‌های ساخته‌شده با عملگرهای تعریف‌شده را باید چنین فهمید:

\begin{tagenumerate}{defTrue,defFalse,defNot,defOr,defAnd,defIf,defIff,defEx,defAll}

\tagitem{defTrue}{$\ltrue$ اختصارِ
  \iftag{prvFalse}{$\lnot\lfalse$}{$(!A \lor \lnot !A)$ برای یک
    !!{formula} اتمی ثابت~$!A$} است.}{}

\tagitem{defFalse}{$\lfalse$ اختصارِ
  \iftag{prvTrue}{$\lnot\ltrue$}{$(!A \land \lnot !A)$ برای یک
    !!{formula} اتمی ثابت~$!A$} است.}{}

\tagitem{defNot}{$\lnot !A$ اختصارِ $!A \lif \lfalse$ است.}{}

\tagitem{defOr}{$!A \lor !B$ اختصارِ
  \iftag{prvAnd}{$\lnot(\lnot !A \land \lnot !B)$}{$\lnot !A \lif
    !B$} است.}{}

\tagitem{defAnd}{$!A \land !B$ اختصارِ
  \iftag{prvOr}{$\lnot(\lnot !A \lor \lnot !B)$}{$\lnot (!A \lif
    \lnot !B)$} است.}{}

\tagitem{defIf}{$!A \lif !B$ اختصارِ
  \iftag{prvOr}{$(\lnot !A \lor !B)$}{$\lnot (!A \land \lnot !B)$} است.}{}

\tagitem{defIff}{$!A \liff !B$ اختصارِ $(!A \lif !B) \land (!B
  \lif !A)$ است.}{}
\end{tagenumerate}
\end{defn}
\end{tagblock}

\begin{defn}[همانی نحوی]
نماد $\ident$ همانی نحوی میان رشته‌های نمادها را بیان می‌کند؛ یعنی
$!A \ident !B$ اگر و تنها اگر $!A$ و $!B$ رشته‌هایی از نمادها باشند که
طول یکسان دارند و در هر جایگاه نماد یکسانی دارند.
\end{defn}

در دو سوی نماد $\ident$ می‌توان رشته‌هایی را قرار داد که با الحاق به دست
آمده‌اند. برای نمونه، $!A \ident (!B \lor !C)$ یعنی رشتهٔ نمادهای~$!A$
همان رشته‌ای است که با الحاق یک پرانتز باز، رشتهٔ $!B$، نماد $\lor$، رشتهٔ
$!C$ و یک پرانتز بسته، به همین ترتیب، به دست می‌آید. اگر چنین باشد، می‌دانیم
نخستین نماد $!A$ یک پرانتز باز است، $!A$ رشتهٔ $!B$ را به‌عنوان زیررشته‌ای
در بر دارد که از نماد دوم آغاز می‌شود، پس از آن زیررشته نماد $\lor$ می‌آید،
و به همین ترتیب.

\end{document}
